\clearpage
\hypertarget{native-481-january10}{}
\begin{center}
\rule{0.55\linewidth}{0.4pt}\\[2pt]
Astronomische Nachrichten. Band XXVII. Nr. 625. S. 13..16. 29. Januar 1848.\\[2pt]
\rule{0.55\linewidth}{0.4pt}
\end{center}
\begin{center}[\textsc{Gauss} an \textsc{Schumacher}.]\quad Göttingen, 10. Januar 1848.\end{center}

Flora habe ich am 5. im Meridian beobachtet.

\begin{center}
\setlength{\tabcolsep}{4pt}
\begin{tabular}{l|r|r|r}
1848 Januar 5 & $9^{\rm u}\ 11^{\rm m}\ 44^{\rm s}8$ & $62^\circ\ 30'\ 18''1$ & $+16^\circ\ 9'\ 45''2$ \\
\end{tabular}
\end{center}

Die Abweichung der \textsc{Goldschmidt}schen Elemente beträgt $+16''4$, $+3''5$. Am 3. fand ich Flora zwar deutlich lichtschwächer als im Dezember des Vorjahres, aber immer noch hell genug, dass sie wohl fast volle Beleuchtung vertragen hätte. Am 5. musste ich die Beleuchtung stark dämpfen. Tatsächlich vertrug der Planet selbst eine geringe Feldbeleuchtung kaum. Das schrieb ich dem Zustand der Luft zu, da auch die Fixsterne sehr unscharfe und unruhige Bilder lieferten. Jedenfalls ist die Beobachtung vom 5. weniger genau als die vom 3.
